% Standalone LaTeX article (two-column). Build from this directory:
%   pdflatex paper_qGG_mjack_v041826_v2F.tex
%   pdflatex paper_qGG_mjack_v041826_v2F.tex   % rerun for references
% Figures: ../figures/ (relative to this file)
% Source blend: paper_qGG_mjack_v041726 + paper_qGG_mjack_v041826

\documentclass[twocolumn,10pt]{article}

%% Encoding and fonts
\usepackage[T1]{fontenc}
\usepackage[utf8]{inputenc}

%% Page geometry
\usepackage[top=1in,bottom=1in,left=0.75in,right=0.75in]{geometry}

%% Math
\usepackage{amsmath}
\usepackage{amssymb}

%% Graphics
\usepackage{graphicx}

%% Tables
\usepackage{booktabs}
\usepackage{array}
\usepackage{tabularx}
\usepackage{dblfloatfix}
\usepackage{placeins}

%% Colors and hyperlinks
\usepackage[table]{xcolor}
\definecolor{linkblue}{rgb}{0.0,0.2,0.7}
\definecolor{tabheader}{RGB}{230,230,227}
\definecolor{tabrowlight}{RGB}{248,248,246}
\definecolor{tabrowdark}{RGB}{240,240,237}
\usepackage[colorlinks=true,linkcolor=linkblue,citecolor=linkblue,urlcolor=linkblue,bookmarks=true]{hyperref}
\hypersetup{%
  pdftitle={Quantum Kernel Link Prediction on Biomedical Knowledge Graphs},
  pdfauthor={Jonathan Beale, Kevin Robinson, Mark Jack},
  pdfkeywords={quantum machine learning, knowledge graphs, link prediction, drug repurposing, Hetionet, quantum kernel, stacking ensemble}%
}

%% Typography
\usepackage{microtype}

%% Captions
\usepackage{caption}
\captionsetup{font=small}
\usepackage{subcaption}

%% References
\usepackage[numbers,sort&compress]{natbib}

%% Lists
\usepackage{enumitem}

%% Section formatting
\usepackage{titlesec}
\titleformat{\section}{\normalfont\large\bfseries}{\thesection}{1em}{}
\titleformat{\subsection}{\normalfont\normalsize\bfseries}{\thesubsection}{1em}{}

%% Abstract spanning both columns
\usepackage{abstract}
\renewcommand{\abstractnamefont}{\normalfont\bfseries}
\renewcommand{\abstracttextfont}{\normalfont\small}

%% URLs
\usepackage{url}

%% ------------------------------------------------------------------
\title{\textbf{Quantum Kernel Link Prediction on Biomedical Knowledge Graphs:}\\
\textbf{Bridging Graph Embeddings and Quantum Feature Spaces for Drug Repurposing}}

\author{%
Jonathan Beale\textsuperscript{\hspace{-0.05em}$\ast$}%
\and Kevin Robinson \and Mark Jack\\
\small Quantum Global Group, 3100 Ray Ferrero Jr Blvd, Davie, Fl 33314\textsuperscript{\hspace{-0.05em}$\dagger$}%
}

\date{arXiv preprint $\cdot$ April 2026}
%% ------------------------------------------------------------------

\begin{document}
\renewcommand{\thefootnote}{\fnsymbol{footnote}}

\twocolumn[
  \maketitle
  \begin{onecolabstract}
Repurposing approved drugs for new indications bypasses the most expensive and time-consuming stages of \textit{de novo} discovery --- preclinical safety, pharmacokinetic, and manufacturing development --- and can deliver new therapies in 3--6 years instead of the customary 10--15 year, \$2--3 billion pipeline. Biomedical knowledge graphs (KGs) aggregate heterogeneous evidence from genomics, pharmacology, pathways, anatomy, and clinical observation, making them a natural substrate for ranking candidate compound--disease hypotheses. We present the first system to feed learned KG vector embeddings into quantum kernel classifiers for biomedical link prediction, targeting \textit{Compound-treats-Disease} (CtD) relationships on the Hetionet KG (47{,}031 entities, 2.25 million edges, 24 relation types). Full-graph RotatE training maps each entity into a 128-dimensional continuous space in which relations such as ``compound binds gene'' act as approximate rotations in complex space, encoding multi-hop biomedical context that purely molecular descriptors cannot reach. Pairs are then represented by concatenation, elementwise difference, and Hadamard product of the two node embeddings and classified by a stacking ensemble whose base learners are tuned classical tree models (Random Forest, Extra Trees) and a quantum support vector classifier (QSVC) built from a fidelity quantum kernel with a 16-qubit Pauli feature map, combined by a logistic-regression meta-learner trained on held-out base-model probabilities.

A central finding is that the Pauli feature map, which entangles qubits through multi-axis $X$, $Y$, and $Z$ rotations rather than ZZ-only diagonal interactions, induces a Hilbert-space geometry whose decision surface is qualitatively different from the axis-aligned, variance-driven splits of tree-based learners; its per-example errors are therefore weakly correlated with those of the classical base models, so the Pauli-QSVC contributes genuinely non-redundant signal that the stacking meta-learner can exploit --- an ensemble-diversity effect in which a weaker but less-correlated quantum base model outperforms a stronger but redundant one in the final score. Clinical validation against ClinicalTrials.gov, however, reveals a \textit{score-validity inversion problem} in embedding-based scoring: the model's highest-scoring hypotheses are not the ones best supported by independent clinical evidence, because graph-embedding similarity rewards structural proximity in the KG (shared neighborhoods, high-degree hubs, transitive co-occurrence) rather than mechanistic plausibility, allowing densely connected but pharmacologically implausible pairs to score highly. To address this, we introduce a ten-feature mechanism-of-action (MoA) module computed directly from Hetionet's multi-relational structure --- binding-target counts, target-overlap Jaccard, shared-pathway gene counts, pharmacologic-class membership, and chemical and disease similarity to known treaters --- that explicitly encodes whether the compound interacts with the biological machinery implicated in the disease, allowing the meta-learner to down-weight structurally close but mechanistically implausible pairs and re-align the model's score surface with clinical validity.
  \end{onecolabstract}
  \vspace{0.5em}
  {\small\noindent\textbf{Keywords:} quantum machine learning; knowledge graph embeddings; link prediction; drug repurposing; Hetionet; quantum kernel methods; quantum support vector classification; stacking ensemble; mechanism of action.\par}
  \vspace{1em}
]

\footnotetext[1]{Correspondence: \texttt{jonathan.beale@quantumglobalgroup.io}}
\footnotetext[2]{Portions of the text in Sections 1 through 11, Bibliography and appendices were drafted and curated with assistance from Claude Code (Anthropic, 2025), utilizing a user-provided Python machine learning codebase to summarize technical implementations. All output was reviewed, verified, and edited by the authors.}
\setcounter{footnote}{2}
\renewcommand{\thefootnote}{\arabic{footnote}}

%% ==================================================================
\section{Introduction}
%% ==================================================================

\subsection{Motivation}

Drug repurposing --- identifying new therapeutic indications for already-approved compounds --- offers a faster, lower-cost, and lower-risk path to clinical deployment than \textit{de novo} drug discovery. Because the candidate molecules already carry established safety, pharmacokinetic, and manufacturing data, repurposing can compress the typical 10--15 year, \$2--3 billion development timeline into a 3--6 year program at a fraction of the cost. Scientifically, it provides a principled way to exploit the pleiotropy of drug targets and the shared biological mechanisms that link apparently distinct diseases, turning incidental observations --- a kinase inhibitor that also modulates autoimmune inflammation, a cardiovascular drug whose secondary effect delays a neurodegenerative cascade --- into systematic search directions. Computational methods that can prioritize plausible compound--disease hypotheses from a combinatorially vast search space are therefore a key enabling technology.

Biomedical knowledge graphs (KGs) integrate heterogeneous evidence across genomics, pharmacology, pathways, anatomy, and clinical observation, and are a natural substrate for this prioritization. The Hetionet knowledge graph \cite{himmelstein2017} encodes 2.25 million biomedical relationships across 47,031 entities --- genes, compounds, diseases, biological processes, pathways, side effects, and anatomical structures --- with the \textit{Compound-treats-Disease} (CtD) relation capturing 755 known drug--disease treatment pairs. Predicting missing CtD links is a well-defined, clinically interpretable link prediction task that reduces to scoring unseen (compound, disease) pairs and surfacing candidates for further investigation.

KG vector embeddings are the dominant computational tool for such scoring. Methods such as TransE, RotatE, ComplEx, and GNN-based approaches project each entity into a continuous latent space in which the relational structure of the graph is approximately preserved --- in RotatE, for example, each relation is represented as a rotation in complex space, so that $\mathbf{h}+\mathbf{r}\approx\mathbf{t}$ captures directed relational semantics. By training these embeddings on the full Hetionet graph rather than on CtD edges alone, the resulting vectors inherit multi-hop biomedical context --- shared binding partners, pathway co-involvement, disease similarity --- that purely molecular descriptors such as fingerprints and QSAR features cannot express. Candidate (compound, disease) pairs are then represented by simple algebraic combinations of their node embeddings (concatenation, elementwise difference, and Hadamard product) and passed to a classifier; this classifier is where our hybrid quantum-classical contribution sits.

Our model is a stacking ensemble whose base learners are tuned classical tree models (Random Forest and Extra Trees) together with a quantum support vector classifier (QSVC) built from a fidelity quantum kernel with a 16-qubit Pauli feature map. A logistic-regression meta-learner is trained on held-out base-model probabilities so that the quantum and classical arms act as complementary experts rather than as competing single models. The Pauli feature map entangles qubits through multi-axis $X$, $Y$, and $Z$ rotations rather than ZZ-only diagonal interactions, inducing a Hilbert-space geometry whose decision surface is qualitatively different from the axis-aligned, variance-driven splits of tree-based learners. Its per-example errors are therefore weakly correlated with those of the classical base models, and the Pauli-QSVC contributes genuinely non-redundant signal that the meta-learner can exploit --- an ensemble-diversity effect in which a weaker but less-correlated base model raises the overall decision quality more than a stronger but redundant one would.

Classical KG embedding methods are mature \cite{mayers2023} and achieve strong baselines on CtD. Quantum machine learning, meanwhile, offers kernel methods that compute similarity in exponentially large Hilbert spaces via parameterized quantum circuits, and quantum kernels have been applied to drug-target interaction prediction \cite{qkdti2025} and ligand-based virtual screening \cite{kruger2023} --- but exclusively on molecular-level features such as fingerprints and QSAR descriptors. No prior work feeds learned KG embedding vectors into quantum feature spaces for link prediction; the two research tracks --- classical KG-based repurposing and quantum molecular drug discovery --- have developed independently, and this paper bridges them.

Even a well-ranked link predictor, however, is only useful if its ordering aligns with biological reality. Clinical validation of our top novel predictions against ClinicalTrials.gov exposes a \textit{score-validity inversion problem} in embedding-based scoring: the model's highest-scoring hypotheses are not the ones best supported by independent clinical evidence. The inversion arises because graph-embedding similarity rewards structural proximity in the KG --- shared neighborhoods, high-degree hubs, and transitive co-occurrence --- rather than mechanistic plausibility, so pairs that sit in dense, well-connected regions can score highly without any pharmacological rationale for treatment. To address this, we introduce a ten-feature mechanism-of-action (MoA) module computed directly from Hetionet's multi-relational structure: compound binding-target count (CbG), disease-associated gene count (DaG), shared targets, target-overlap Jaccard index, shared-pathway gene count (GpPW), pharmacologic-class membership (PCiC), chemically similar compounds (CrC), similar compounds that are known treaters, similar diseases (DrD), and similar diseases that are known to be treated. These features explicitly encode the question the embedding cannot answer --- \textit{does the compound interact with the biological machinery implicated in the disease?} --- and, when concatenated with the embedding-based features, allow the meta-learner to down-weight structurally close but mechanistically implausible pairs and re-align the model's score surface with clinical validity.

\subsection{Problem Statement}

Given Hetionet with known CtD edges as positives and hard-sampled negatives, train a binary classifier predicting whether an unseen (compound, disease) pair represents a treatment relationship. We compare classical-only, quantum-only, and hybrid quantum-classical stacking ensemble approaches, reporting test PR-AUC as the primary metric --- appropriate for the class-imbalanced link prediction setting.

\subsection{Contributions}

\begin{enumerate}[leftmargin=*, label=\arabic*.]

\item \textbf{Novel intersection.} To our knowledge, this is the first system to combine learned KG embedding vectors (RotatE, full-graph) with quantum kernel classifiers (QSVC, fidelity kernel) for biomedical link prediction, bridging two previously disjoint research tracks.

\item \textbf{Counterintuitive quantum-classical synergy.} The Pauli feature map degrades standalone QSVC performance (0.7216 $\to$ 0.6343, $-8.7$ pp) relative to ZZ, yet raises ensemble PR-AUC (0.7408 $\to$ 0.7987, $+5.8$ pp). The Pauli kernel generates predictions with lower correlation to classical tree model errors, providing greater diversity for the stacking meta-learner.

\item \textbf{Mechanism-of-action feature module.} We introduce 10 pharmacological plausibility features derived from multi-relational Hetionet structure --- binding target overlap (CbG $\cap$ DaG), pathway co-involvement (GpPW), pharmacologic class membership (PCiC), and chemical and disease similarity (CrC, DrD) --- to address the score-validity inversion problem inherent in embedding-based scoring.

\item \textbf{Clinical validation.} Top predictions are validated against ClinicalTrials.gov. Two of six novel predictions (Losartan$\to$atherosclerosis, Mitomycin$\to$liver cancer) are supported by 7+ clinical trials each; one (Ezetimibe$\to$gout) represents a biologically plausible novel hypothesis.

\item \textbf{Reproducible pipeline} with configurable embeddings (RotatE, ComplEx, DistMult), feature maps (ZZ, Pauli), dimensionality reduction strategies, ensemble methods, GPU-accelerated simulation, and IBM Quantum hardware backends.

\end{enumerate}

%% ==================================================================
\section{Background and Related Work}
%% ==================================================================

\subsection{Hetionet and KG-Based Drug Repurposing}

Hetionet v1.0 \cite{himmelstein2017} is a heterogeneous biomedical knowledge graph with 47,031 nodes and 2,250,198 edges across 24 relation types. Key relations include: Compound-treats-Disease (CtD, 755 edges), Compound-binds-Gene (CbG, 11,571), Disease-associates-Gene (DaG, 12,623), Gene-participates-Pathway (GpPW, 84,372), Compound-resembles-Compound (CrC, 6,486), and Disease-resembles-Disease (DrD, 543). The entity set covers 1,552 compounds, 137 diseases, and 20,945 genes, among others.

Mayers et al.\ \cite{mayers2023} benchmarked seven KG embedding methods on biomedical KGs, including RotatE and ComplEx, achieving MRR of 0.9792 via ensemble. Recent work has extended KG-based repurposing with LLM-assisted retrieval augmentation \cite{llmrag2025} and multi-database KG construction for disease-specific targets \cite{quantum2026}. All are exclusively classical approaches operating over embedding scores or graph traversal.

\subsection{Quantum Kernel Methods}

Quantum kernels compute pairwise similarity in a quantum feature space: $k(\mathbf{x}, \mathbf{y}) = |\langle 0 | U^\dagger(\mathbf{x})\, U(\mathbf{y}) | 0 \rangle|^2$, where $U(\mathbf{x})$ is a parameterized encoding circuit (feature map). A quantum support vector classifier (QSVC) uses this kernel as a drop-in replacement for classical kernels. The expressivity of the kernel depends on the feature map architecture --- ZZ (second-order Pauli-Z interactions) or Pauli (mixed X, Y, Z, ZZ interactions) --- and the number of circuit repetitions. Kernel computation scales as $O(n^2)$ in training set size, placing practical constraints on dataset size without approximation \cite{schuld2021}.

\subsection{Quantum Methods in Drug Discovery}

QKDTI \cite{qkdti2025} applied Quantum Support Vector Regression with quantum feature mapping to drug-target interaction prediction using molecular descriptors. Kruger et al.\ \cite{kruger2023} demonstrated QSVC for ligand-based virtual screening on molecular fingerprints. A 2024 preprint \cite{hybrid2024} described a hybrid quantum-classical pipeline for real-world drug discovery, again at the molecular level. A 2026 preprint \cite{quantum2026} used a Coherent Ising Machine for allosteric site detection and molecular docking.

\textbf{The critical distinction:} all prior quantum drug discovery work applies quantum kernels or circuits to molecular-level features --- fingerprints, QSAR descriptors, binding affinities. All prior KG-based drug repurposing uses classical methods. Our work is the first to apply quantum kernels to \textit{learned KG embedding vectors}, encoding the relational structure of the full biomedical knowledge graph as quantum feature space inputs.

\subsection{Stacking Ensembles and Model Diversity}

Stacking trains a meta-learner on base model predictions, exploiting the principle that diverse errors --- not individually high performance --- drive ensemble gains \cite{wolpert1992}. Hybrid quantum-classical ensembles extend this pattern by adding a quantum base learner that computes similarity in a reproducing-kernel Hilbert space inaccessible to efficient classical simulation in the large-circuit limit \cite{havlicek2019,schuld2019,liu2021,abbas2021,incudini2023}. We exploit this explicitly: the Pauli QSVC performs below the ZZ QSVC in isolation but provides more diverse errors relative to random forest and extra trees, enabling the meta-learner to extract greater combined signal (Section~\ref{sec:hybrid}).

%% ==================================================================
\section{Dataset}
%% ==================================================================

\begin{table*}[t]
\centering
\caption{Hetionet CtD Dataset Statistics}
\label{tab:dataset}
\begin{tabular}{@{}ll@{}}
\toprule
\textbf{Statistic} & \textbf{Value} \\
\midrule
Source & Hetionet v1.0 (het.io) \\
Total entities & 47,031 \\
\quad-- Compounds & 1,552 \\
\quad-- Diseases & 137 \\
\quad-- Genes & 20,945 \\
\quad-- Other (pathways, etc.) & 24,597 \\
Total edges (all 24 relations) & 2,250,198 \\
Target relation: CtD & 755 positive edges \\
Train positives (80\%) & 604 \\
Test positives (20\%) & 151 \\
Negative sampling strategy & Hard negatives, 1:1 ratio \\
Train pairs (pos.\ + hard neg.) & 1,208 \\
Test pairs & 302 \\
Feature dim.\ --- classical path & 1177 (base) / 1187 (+MoA) \\
Feature dim.\ --- quantum path & 16 qubits (1177D$\to$24D$\to$16Q) \\
Class balance (constructed) & 50 / 50 \\
Unique compounds in train heads & 354 (29.3\% of 1,208 pairs) \\
Unique diseases in train tails & 75 (6.2\% of 1,208 pairs) \\
\bottomrule
\end{tabular}
\end{table*}

The low tail uniqueness (75 unique diseases from 137 total) reflects that Hetionet's CtD subgraph is concentrated: most training pairs map compounds to a small set of well-studied diseases. This is a structural property of the graph, not a data quality issue, and is relevant to understanding QSVC kernel informativeness (see Section~\ref{sec:results_feature_map}).

\FloatBarrier
%% ==================================================================
\section[Hybrid pipeline for CtD link prediction]{Hybrid quantum-classical pipeline for Compound-treats-Disease (CtD) link prediction on Hetionet}
%% ==================================================================

\subsection{Pipeline Overview}
\label{sec:pipeline}

The full pipeline proceeds as follows:

\begin{enumerate}[leftmargin=*, label=\arabic*.]
\item Load Hetionet and extract CtD edges as positive examples.
\item Train full-graph RotatE embeddings across all 24 relation types using PyKEEN \cite{pykeen2021}.
\item Construct train/test pair sets with hard negative sampling at 1:1 ratio.
\item Build feature vectors for each (compound, disease) pair from entity embeddings and graph topology.
\item \textbf{Classical path:} Train RandomForest, ExtraTrees, and LogisticRegression with optional GridSearchCV tuning.
\item \textbf{Quantum path:} Reduce dimensionality via PCA, encode into a Pauli or ZZ feature map, compute fidelity quantum kernel, train QSVC.
\item \textbf{Stacking ensemble:} Train a logistic regression meta-learner on out-of-fold base model predictions.
\item Evaluate on held-out test set; report PR-AUC and ROC-AUC.
\end{enumerate}

\begin{figure*}[t]
\centering
\includegraphics[width=\textwidth]{../figures/fig1_pipeline.pdf}
\caption{Hybrid quantum-classical pipeline for Compound-treats-Disease (CtD) link prediction on Hetionet. Full-graph RotatE embeddings are trained over all 24 relation types (2.25M edges). Pair-wise features are constructed from entity embeddings and training-graph topology. The quantum path applies a 16-qubit Pauli feature map to a PCA-reduced representation; the classical path uses RF/ET/LR with GridSearchCV tuning. A logistic regression stacking meta-learner combines base model outputs.}
\label{fig:pipeline}
\end{figure*}

\subsection{Software architecture and orchestration}

The reference implementation is orchestrated by \texttt{scripts/run\_optimized\_pipeline.py}. It composes: \textbf{\texttt{kg\_layer}} --- Hetionet loading, CtD extraction, hard-negative generation (\texttt{kg\_loader}), train-only graph construction for topology features, \texttt{EnhancedFeatureBuilder} for embedding and structural features, and optional MoA computation (\texttt{moa\_features}); \textbf{\texttt{quantum\_layer}} --- PCA and linear projection to qubits, Qiskit feature maps, fidelity-kernel QSVC; and \textbf{classical baselines} --- scikit-learn RandomForest, ExtraTrees, LogisticRegression, with optional GridSearchCV, plus stacking. End-to-end behavior is controlled by the same flags as the reproducibility command (Section~\ref{sec:exp_repro}).

\subsection{Supervised pairs, hard negatives, and leakage controls}

\textbf{Train/test split.} Known CtD positives are split 80\%/20\% into training and test sets (random seed~42). For each side we sample an equal number of negatives so constructed train and test sets are balanced 1:1 positive/negative.

\textbf{Hard negatives.} With \texttt{--negative\_sampling hard}, negatives are generated by \texttt{get\_hard\_negatives} in \texttt{kg\_loader} using the default \textit{degree\_corrupt} strategy: for each positive $(s,t)$, either the head or tail is replaced (Bernoulli with fixed tail-corruption probability) by a \emph{degree-weighted} alternative entity of the same role, rejecting pairs that coincide with a known positive edge. This yields negatives that are structurally closer to real positives than uniform random sampling. If too few unique negatives are obtained within a budget, the pipeline tops up with random negatives.

\textbf{Leakage controls.} The NetworkX graph used for neighborhood statistics, and the edge list passed into domain features, include \textbf{training} positives and their negatives only --- not test positives. MoA features that reference known treatments use \textbf{training} CtD edges exclusively (Section~\ref{sec:moa}). \texttt{StandardScaler} is fit on training features and applied to test features without refitting.

\subsection{Full-Graph Knowledge Graph Embeddings}

We train RotatE \cite{rotate2019} embeddings over all 2,250,198 Hetionet edges across all 24 relation types using PyKEEN \cite{pykeen2021}. RotatE models each relation as an element-wise rotation in complex space: for a valid triple $(h, r, t)$, the relation $r$ acts as
\begin{equation}
\mathbf{h} \circ \mathbf{r} \approx \mathbf{t}
\end{equation}
where $\circ$ denotes complex-space element-wise multiplication. Training on the full graph --- rather than only the 755 CtD edges --- enriches compound and disease entity representations with signal from gene binding (CbG), disease-gene associations (DaG), pathway participation (GpPW), side effects (CcSE), chemical similarity (CrC), and all other relation types.

\textbf{Primary configuration:} 128D, 200 epochs, full-graph mode (\texttt{--full\_graph\_embeddings}), hard negatives for embedding training, random seed 42.

\textbf{Extended configuration:} 256D, 250 epochs, full-graph mode.

\subsection{Pair feature construction}

For each (compound $c$, disease $d$) pair, \texttt{EnhancedFeatureBuilder} assembles a vector $\mathbf{x}_{cd}$ from three blocks (graph and domain features on; MoA off unless \texttt{--use\_moa\_features}). Let $\mathbf{h}=\mathbf{e}_c$ and $\mathbf{t}=\mathbf{e}_d$ denote head and tail embeddings in $\mathbb{R}^{d}$ (primary run: $d=128$).

\textbf{(i) Embedding block $\phi(\mathbf{h},\mathbf{t})$.} We concatenate $\mathbf{h}$ and $\mathbf{t}$; elementwise $|\mathbf{h}-\mathbf{t}|$ and $\mathbf{h}\odot\mathbf{t}$; three global similarity scalars (cosine similarity, Euclidean distance, dot product); $(\mathbf{h}+\mathbf{t})/2$; elementwise $\mathbf{h}^{2}$, $\mathbf{t}^{2}$, $\sqrt{|\mathbf{h}|}$, and $\sqrt{|\mathbf{t}|}$. This yields dimension $9d+3$ (1155 when $d=128$).

\textbf{(ii) Graph topology $\mathbf{g}_{cd}$ (14 dimensions).} From a NetworkX graph built on \textbf{training} edges only: head/tail degree, common-neighbor count, Jaccard coefficient, Adamic--Adar, preferential attachment, resource allocation, shortest-path length (or placeholder if disconnected), PageRank (head, tail), betweenness (head, tail), clustering (head, tail).

\textbf{(iii) Domain features $\mathbf{d}_{cd}$ (8 dimensions).} Binary entity-type indicators for head and tail (compound / disease / gene) and counts of distinct metaedge types incident to each endpoint in the training edge list.

Thus $\mathbf{x}_{cd} = [\phi(\mathbf{h},\mathbf{t})\,\|\,\mathbf{g}_{cd}\,\|\,\mathbf{d}_{cd}]$ has dimension $1177$ for $d{=}128$. The optional MoA block adds 10 dimensions (Section~\ref{sec:moa}). The extended 256D embedding run uses the same recipe with $d{=}256$, giving a larger classical vector before PCA.

\subsection{Classical Base Models}
\label{sec:classical}

The classical arm of the pipeline consists of three base learners operating directly on the 1177-dimensional pair feature vector $\mathbf{x}_{cd}$:

\begin{itemize}[leftmargin=*]
\item \textbf{Random Forest (RF).} 500 estimators in full runs (reduced estimator counts when \texttt{--fast\_mode} is enabled), Gini split criterion; \texttt{max\_depth}, \texttt{min\_samples\_split}, \texttt{min\_samples\_leaf}, and \texttt{max\_features} tuned by 5-fold GridSearchCV when \texttt{--tune\_classical} is set.
\item \textbf{Extra Trees (ET).} 500 estimators with fully randomised split thresholds, same hyperparameter grid as RF.
\item \textbf{Logistic Regression (LR).} L2-regularised logistic regression with inverse-regularisation strength $C_{\mathrm{LR}}$ selected over $[10^{-3}, 10^{2}]$ log-uniform; the LR baseline also serves as a linear reference for the more expressive tree and quantum learners and, at the meta level, as the combiner of the stacking ensemble.
\end{itemize}

All three produce calibrated probability outputs $p_{m}(y=1|\mathbf{x}_{cd})$ that are exposed to the stacking meta-learner of Section~\ref{sec:hybrid}. RF and ET are retained for their robustness to noisy and redundant features in the high-dimensional embedding representation, LR for its interpretability and low variance. Their piecewise-constant (RF, ET) and linear (LR) decision surfaces will be shown, in Section~\ref{sec:hybrid}, to be complementary to the continuous, globally-entangled decision surface of the quantum kernel classifier.

\subsection{Quantum Kernel Classifier}
\label{sec:qsvc}

\textbf{Dimensionality reduction.} The $1177$-dimensional classical feature vector is compressed to 24D via PCA, then projected to 16 qubits by a learned linear layer. For the 256D embedding configuration, no pre-PCA is applied and projection goes directly to 12 qubits. This reduction is the primary information bottleneck in the quantum path.

\textbf{Feature maps.} We evaluate two Qiskit \cite{qiskit2023} feature maps:
\begin{itemize}[leftmargin=*]
\item \textbf{ZZ feature map:} Second-order Pauli-Z interactions between adjacent qubits; reps=2--3. Equivalent to the \textit{ZZFeatureMap} in Qiskit's circuit library.
\item \textbf{Pauli feature map:} Mixed Pauli interactions (X, Y, Z, ZZ) over all qubit pairs; reps=2. Equivalent to \textit{PauliFeatureMap} with full entanglement. More expressive than ZZ but produces a different kernel geometry.
\end{itemize}

\textbf{Fidelity quantum kernel.} The quantum kernel is:

\begin{equation}
k_{Q}(\mathbf{x}, \mathbf{y}) = \left|\langle 0^{\otimes n} | U^\dagger(\mathbf{x})\, U(\mathbf{y}) | 0^{\otimes n} \rangle\right|^2
\label{eq:qkernel}
\end{equation}

computed via statevector simulation. No Nystr\"{o}m approximation is used in the primary run (\texttt{nystrom\_m=None}), resulting in a full $1208 \times 1208$ kernel matrix requiring approximately 1.46 million circuit evaluations.

\textbf{Computation time.} The primary 16-qubit Pauli kernel computation required \textbf{2,619 seconds} (43.6 minutes) on CPU statevector simulator --- confirming this is a genuine full-dataset quantum kernel, not a cached or subsampled approximation. The extended 12-qubit configuration computed in approximately 99 seconds (first trial; subsequently cached for Optuna sweeps).

\textbf{QSVC.} Regularization $C_{Q}=0.1$ (primary, manually set) or $C_{Q}=0.676$ (extended, Optuna-tuned). Classification is via scikit-learn's SVC with the precomputed quantum kernel matrix, yielding a standalone probability $p_{\mathrm{QSVC}}(y=1|\mathbf{x}_{cd})$ via Platt scaling on the cross-validated fold split.

\subsection{Hybrid classification model: stacking ensemble of quantum and classical meta-learners}
\label{sec:hybrid}

\subsubsection*{Related work on hybrid ensemble classification}

Hybrid quantum-classical ensembles have emerged as a pragmatic answer to a fundamental limitation of near-term quantum hardware: quantum models are often individually weaker than their classical counterparts, yet they compute similarity in reproducing-kernel Hilbert spaces whose geometry is provably inaccessible to efficient classical simulation in the large-circuit limit \cite{havlicek2019,liu2021,schuld2019}. Combining the two through a classical meta-learner therefore offers a route to accuracy improvements that neither arm can achieve alone. The algorithmic backbone is Wolpert's stacked generalisation \cite{wolpert1992}: train a set of base classifiers, generate held-out predictions from each, and train a second-level meta-learner on those predictions. Schuld and Killoran \cite{schuld2019} showed that quantum feature maps induce reproducing-kernel Hilbert spaces, making quantum support vector classifiers natural stacking candidates, while Havlic\v{c}ek et al.\ \cite{havlicek2019} and Schuld and Petruccione \cite{schuld2021} established the fidelity quantum kernel as the canonical kernel for supervised learning in those spaces. Abbas et al.\ \cite{abbas2021} demonstrated that the effective dimension of quantum models can exceed that of comparable classical neural networks, which, in an ensembling context, translates into complementary rather than redundant capacity. Several recent studies have explicitly combined quantum and classical classifiers in a stacked or voted ensemble for particle-physics event tagging, medical image classification, and tabular fraud detection \cite{incudini2023}; a recurring empirical observation in that literature --- and the one our own experiments confirm --- is that the ensemble gain is driven not by the raw accuracy of the quantum arm but by the \emph{diversity} of its errors relative to the classical arm. Our contribution is to extend this pattern to biomedical link prediction on a multi-relational knowledge graph, where both ensemble arms operate on the same 128-dimensional RotatE embedding but resolve it through fundamentally different geometric lenses.

\subsubsection*{Algorithmic setup of the stacking ensemble}

The ensemble is organised as a two-level hierarchy. At level zero sit the four base learners defined in Sections~\ref{sec:classical} and~\ref{sec:qsvc}: Random Forest, Extra Trees, Logistic Regression, and the quantum support vector classifier. Each is trained on the full training set to produce a class-probability output $p_{m}(y=1|\mathbf{x}_{cd})$ for $m \in \{\mathrm{RF},\mathrm{ET},\mathrm{LR},\mathrm{QSVC}\}$. To prevent the meta-learner from seeing base predictions over-fit to the same examples it will later combine, we use a $K=5$ stratified cross-validation protocol to generate \emph{out-of-fold} (OOF) probabilities: for each fold $k \in \{1,\dots,K\}$, each base model is trained on folds $\{1,\dots,K\}\setminus\{k\}$ and used to predict on the held-out fold. The stacked probability vector

\begin{equation}
\mathbf{z}_i = \bigl[\,p_{\mathrm{RF}}^{\mathrm{OOF}}(\mathbf{x}_i),\, p_{\mathrm{ET}}^{\mathrm{OOF}}(\mathbf{x}_i),\, p_{\mathrm{LR}}^{\mathrm{OOF}}(\mathbf{x}_i),\, p_{\mathrm{QSVC}}^{\mathrm{OOF}}(\mathbf{x}_i)\,\bigr] \in [0,1]^{4}
\label{eq:stack}
\end{equation}

is therefore uncontaminated by in-sample leakage. At level one, a logistic-regression meta-learner with L2 penalty learns the weights $\mathbf{w} \in \mathbb{R}^{4}$ and bias $b \in \mathbb{R}$ that map $\mathbf{z}_i$ to the final probability $\hat{p}(\mathbf{x}_i) = \sigma(\mathbf{w}^{\top}\mathbf{z}_i + b)$, where $\sigma$ is the logistic sigmoid. After the meta-learner is trained, each base model is re-fit on the \emph{full} training set to maximise its standalone calibration, and the four meta-level weights are held fixed at inference. The resulting scoring function is a single, differentiable composition that blends quantum and classical evidence through a learned linear layer over base-model probabilities.

\subsubsection*{The quantum arm as a meta-learner}

The quantum arm contributes to the meta-space via the Platt-scaled QSVC probability derived from the fidelity kernel of Eq.~\eqref{eq:qkernel}. Although a standalone QSVC is not typically thought of as a ``meta-learner,'' its role inside the stacking ensemble is precisely that: its output is not consumed by an end user directly but is re-interpreted by a higher-level linear model as one of several expert opinions. Two properties of the quantum arm make it well-suited to this role. First, the fidelity kernel $k_{Q}$ is a continuous, symmetric, positive-semidefinite function whose level sets are defined by interference patterns in a $2^{n}$-dimensional Hilbert space; the decision surface is therefore globally entangled and does not align with any coordinate axis of the input feature vector. Second, the Pauli feature map --- which encodes $\mathbf{x}$ through multi-axis $X$, $Y$, and $Z$ rotations with full-graph entanglement --- produces a kernel that is qualitatively distinct from the ZZ-only diagonal feature map, yielding predictions that are weakly correlated with those of the classical base learners (Section~\ref{sec:results_feature_map}). This weak correlation is what gives the quantum arm its value as a meta-input.

\subsubsection*{The classical arm as a meta-learner}

Logistic regression plays a dual role in our architecture: at level zero as a base learner with decision function

\begin{equation}
p_{\mathrm{LR}}(y=1|\mathbf{x}) = \sigma\bigl(\mathbf{w}_{\mathrm{LR}}^{\top}\mathbf{x} + b_{\mathrm{LR}}\bigr),
\end{equation}

trained by minimising the regularised cross-entropy

\begin{equation}
\mathcal{L}(\mathbf{w}_{\mathrm{LR}}, b_{\mathrm{LR}}) = -\frac{1}{n}\sum_{i=1}^{n}\bigl[y_i \log p_i + (1-y_i)\log(1-p_i)\bigr] + \lambda \|\mathbf{w}_{\mathrm{LR}}\|_{2}^{2},
\end{equation}

and at level one as the meta-combiner of Eq.~\eqref{eq:stack}. Using the same functional form at both levels makes the learned meta-weights interpretable as evidence strengths --- each base model contributes a coefficient $w_{m}$ that, in absolute value, measures how much the meta-learner trusts expert $m$ after cross-validated calibration. Random Forest and Extra Trees complete the classical arm, each training a bagged ensemble of axis-aligned decision trees on bootstrapped subsets of the training data. Their decision boundaries are piecewise-constant partitions of the 1177-dimensional embedding-derived feature space, which stands in sharp contrast to the continuous fidelity-kernel surface of the QSVC. The three classical base learners therefore span the local (trees) and global-linear (LR) regimes of the feature space, and the QSVC adds a global-entangled regime that neither can reach.

\subsubsection*{Why the collaboration beats isolated classifiers}

The advantage of the hybrid ensemble over any single classifier can be made algebraically transparent. For base models with per-example errors $\epsilon_{m}(\mathbf{x}) = p_{m}(\mathbf{x}) - y$, the variance of the weighted meta-probability $\hat{p}(\mathbf{x}) = \sum_{m} w_{m} p_{m}(\mathbf{x})$ is

\begin{equation}
\mathrm{Var}\bigl[\hat{p}(\mathbf{x})\bigr] = \sum_{m,m'} w_{m} w_{m'}\,\mathrm{Cov}\bigl(\epsilon_{m}(\mathbf{x}), \epsilon_{m'}(\mathbf{x})\bigr).
\label{eq:ensemble_var}
\end{equation}

When two base models produce weakly correlated errors, the off-diagonal covariance terms in Eq.~\eqref{eq:ensemble_var} shrink and the meta-prediction's variance is reduced at fixed bias. Because the Pauli-QSVC errors are weakly correlated with those of RF, ET, and LR --- empirically quantified in Section~\ref{sec:results_feature_map} --- the ensemble achieves lower variance at the decision boundary than any single base model, which translates into the PR-AUC gain of Table~\ref{tab:primary}. An equally intuitive geometric statement is: the classical arm supplies sharp, axis-aligned decisions in regions of the feature space where the embedding coordinates are directly informative; the quantum arm supplies continuous, globally-entangled similarity in regions where coordinate-aligned geometry alone is misleading. The meta-learner has simply to assign non-negative, appropriately scaled weights to each signal. A notable corollary, first documented in our ZZ-vs-Pauli ablation, is that a quantum base learner whose standalone PR-AUC is \emph{lower} can nevertheless produce a \emph{higher-performing} ensemble if its errors are more decorrelated from those of the classical arm; this reframes the optimisation target for the quantum component from standalone accuracy to classical-quantum error independence.

\subsubsection*{Parameter set of the ensemble}

The end-to-end hybrid classifier exposes a hyperparameter set
\begin{equation}
\boldsymbol{\theta} = \bigl(\boldsymbol{\theta}_{\mathrm{RF}},\, \boldsymbol{\theta}_{\mathrm{ET}},\, \boldsymbol{\theta}_{\mathrm{LR}},\, \boldsymbol{\theta}_{\mathrm{QSVC}},\, \boldsymbol{\theta}_{\mathrm{meta}},\, \boldsymbol{\theta}_{\mathrm{feat}}\bigr),
\end{equation}
with components:
\begin{itemize}[leftmargin=*]
\item $\boldsymbol{\theta}_{\mathrm{RF}} = (n_{\mathrm{estimators}},\, \mathrm{max\_depth},\, \mathrm{min\_samples\_split},\, \mathrm{min\_samples\_leaf},\, \mathrm{max\_features})$, with $n_{\mathrm{estimators}} \in \{200, 500, 1000\}$, $\mathrm{max\_depth} \in \{\text{None}, 10, 20, 30\}$, $\mathrm{min\_samples\_split} \in \{2, 5, 10\}$.
\item $\boldsymbol{\theta}_{\mathrm{ET}}$ follows the same grid.
\item $\boldsymbol{\theta}_{\mathrm{LR}} = (C_{\mathrm{LR}},\, \text{penalty})$ with $C_{\mathrm{LR}} \in [10^{-3}, 10^{2}]$ log-uniform.
\item $\boldsymbol{\theta}_{\mathrm{QSVC}} = (C_{Q},\, \text{feature\_map},\, \text{reps},\, n_{\text{qubits}},\, \text{pre\_pca\_dim})$ with $C_{Q} \in [10^{-2}, 10^{1}]$ log-uniform, $\text{feature\_map} \in \{\mathrm{ZZ}, \mathrm{Pauli}\}$, $\text{reps} \in \{1, 2, 3\}$, $n_{\text{qubits}} \in \{8, 12, 16\}$, $\text{pre\_pca\_dim} \in \{0, 16, 24, 32\}$.
\item $\boldsymbol{\theta}_{\mathrm{meta}} = (C_{\mathrm{meta}},\, \text{penalty},\, \text{class\_weight})$ for the level-1 logistic regression.
\item $\boldsymbol{\theta}_{\mathrm{feat}}$ toggles the MoA feature module (Section~\ref{sec:moa}) and the embedding construction variant.
\end{itemize}
Together these parameters span a mixed continuous/discrete search space of approximately $2 \times 10^{7}$ configurations, which rules out exhaustive grid search.

\subsubsection*{Optuna-driven hyperparameter optimisation}

To navigate this space efficiently we use Optuna \cite{optuna2019}, a sequential model-based optimisation library that maintains a Tree-structured Parzen Estimator (TPE) surrogate of the objective as it observes trial outcomes. The objective is the mean five-fold cross-validated PR-AUC of the stacking ensemble on the training set,
\begin{equation}
J(\boldsymbol{\theta}) = \frac{1}{K}\sum_{k=1}^{K}\mathrm{PR\text{-}AUC}\bigl(\hat{p}^{(-k)}_{\boldsymbol{\theta}},\, \mathbf{y}^{(k)}\bigr),
\label{eq:optuna_obj}
\end{equation}
where $\hat{p}^{(-k)}_{\boldsymbol{\theta}}$ is the ensemble prediction on fold $k$ trained on the remaining folds with hyperparameters $\boldsymbol{\theta}$. Optuna's TPE sampler proposes the next $\boldsymbol{\theta}^{(t)}$ to maximise an expected-improvement acquisition against the running best $J^{\star}$, while its Asynchronous Successive Halving pruner terminates unpromising trials early by comparing intermediate fold scores against a quantile of previously observed trials.

A practical complication is that $\boldsymbol{\theta}_{\mathrm{QSVC}}$ controls the quantum kernel itself, whose computation dominates runtime (2{,}619\,s for the 16-qubit Pauli configuration on CPU statevector simulation). Na\"{i}ve Optuna over $\boldsymbol{\theta}$ would require a kernel recomputation at every trial. We therefore layer the study in two stages. An \emph{outer quantum-structure study} (20 trials) varies $(\text{feature\_map}, \text{reps}, n_{\text{qubits}}, \text{pre\_pca\_dim})$ while holding the classical hyperparameters at reasonable defaults; for each distinct structural tuple the kernel matrix $K_{Q}$ is computed once and cached to disk keyed on $(\text{feature\_map}, \text{reps}, n_{\text{qubits}}, \text{pre\_pca\_dim}, \mathrm{seed})$. An \emph{inner classical-hyperparameter study} (60--80 trials) then varies $C_{Q}$ together with the classical and meta parameters on a frozen kernel matrix, requiring no further circuit evaluations. This two-stage design reduces the overall Optuna budget from a nominal $O(10^{3})$ quantum kernel computations to 5--10 kernel computations --- the only regime in which Bayesian optimisation of the full $\boldsymbol{\theta}$ is tractable on current simulators.

The optimiser is seeded with $\mathtt{seed}=42$ for both the TPE sampler and each inner trial's train/test split, so that reported Optuna-tuned results are deterministically reproducible. Convergence is monitored through the study's \texttt{best\_value} trajectory: in the extended configuration (Table~\ref{tab:extended}) the study was judged converged after 83 classical trials had failed to improve the running best by more than $10^{-3}$ PR-AUC over a rolling 20-trial window. The final hyperparameter point $\boldsymbol{\theta}^{\star}$ is written to the JSON result file alongside model metrics, so every reported number in this paper is traceable to the exact study configuration that produced it.

%% ==================================================================
\section{Mechanism-of-Action Feature Module}
\label{sec:moa}
%% ==================================================================

\subsection{Motivation: structural proximity vs.\ mechanistic plausibility}

Graph-embedding scores measure \emph{structural proximity} in the knowledge graph. Structural proximity correlates with treatment plausibility on average --- drugs that treat a disease tend to bind genes associated with it, share pathways with its causal genes, and resemble other compounds known to treat the same or related conditions --- but the correlation is far from perfect. In the high-density, high-degree regions of Hetionet, a compound and a disease can be close in embedding space simply because both are connected to the same hub genes or to related comorbidities, with no pharmacological basis for treatment. The Abacavir~$\to$~ocular cancer case (Section~\ref{sec:clinical}) is a clean example: HIV comorbidity pulls Abacavir and ocular malignancies into the same neighbourhood of the graph even though Abacavir's mechanism of action --- reverse-transcriptase inhibition --- has nothing to do with ocular-cancer biology.

The mechanism-of-action (MoA) feature module is designed to give the classifier an orthogonal signal that reflects \emph{mechanistic plausibility} rather than structural proximity. Ten features are computed directly from the multi-relational Hetionet subgraph for every (compound, disease) pair. They are organised into three biological channels --- direct target evidence, analogical drug-class evidence, and analogical disease-class evidence --- each with a clear pharmacological interpretation.

\subsection{Feature definitions}

Let $\mathcal{T}_{\mathrm{train}}$ denote the set of (compound, disease) pairs with a known CtD edge in the training set. The ten features are defined in Table~\ref{tab:moa} with the signal they capture.

\begin{table*}[t]
\centering
\caption{Mechanism-of-Action Feature Definitions}
\label{tab:moa}
\small
\setlength{\tabcolsep}{4pt}
\begin{tabularx}{\textwidth}{@{}c >{\ttfamily\raggedright\arraybackslash}p{0.22\textwidth} >{\raggedright\arraybackslash}p{0.22\textwidth} X@{}}
\toprule
\textbf{\#} & \textbf{Feature Name} & \textbf{Source Relations / Formula} & \textbf{Signal Captured} \\
\midrule
1 & binding\_targets & $|T(c)|$, CbG & Number of genes the compound binds (graph coverage for $c$). \\
2 & disease\_genes & $|G(d)|$, DaG & Number of genes associated with the disease (graph coverage for $d$). \\
3 & shared\_targets & $|T(c)\cap G(d)|$ & \textbf{Genes bound by $c$ AND linked to $d$ --- direct mechanistic evidence.} \\
4 & target\_overlap & $|T(c)\cap G(d)| / |T(c)\cup G(d)|$ & Jaccard-normalised mechanistic overlap. \\
5 & shared\_pathway\_genes & $|\{g\in T(c): P(g)\cap P(G(d))\ne\emptyset\}|$ & Targets of $c$ that share a pathway with disease genes. \\
6 & pharmacologic\_classes & $|\mathrm{PC}(c)|$, PCiC & Number of pharmacologic classes $c$ belongs to. \\
7 & compound\_similarity & $|S(c)|$, CrC & Chemical-neighbourhood size of $c$. \\
8 & similar\_compounds\_treat & $|\{c'\in S(c): \exists d',\, (c',d')\in\mathcal{T}_{\mathrm{train}}\}|$ & \textbf{Chemically similar compounds that are known treaters --- analogical evidence.} \\
9 & disease\_similarity & $|S(d)|$, DrD & Phenotypic-neighbourhood size of $d$. \\
10 & similar\_diseases\_treated & $|\{d'\in S(d): \exists c',\, (c',d')\in\mathcal{T}_{\mathrm{train}}\}|$ & \textbf{Similar diseases known to be treated --- analogical evidence.} \\
\bottomrule
\end{tabularx}
\end{table*}

\subsection{Feature selection rationale}

Selection was guided by three constraints. First, every feature must be \emph{computable from public Hetionet data alone}, with no reliance on proprietary pharmacology databases; this keeps the module reproducible and re-usable on any multi-relational biomedical KG that shares the same relation vocabulary (CbG, DaG, GpPW, PCiC, CrC, DrD). Second, every feature must be \emph{semantically interpretable} as evidence a pharmacologist would accept: direct binding of a drug target implicated in the disease (features~3--5), membership of a drug class with established activity against the indication class (feature~6), or analogy to known treatments via chemical or disease similarity (features~7--10). Third, the features must be as \emph{orthogonal to RotatE embedding geometry} as possible. Whereas RotatE learns smooth, continuous coordinates from triples, the MoA features are explicit, discrete set-operation counts over named biological relations; this difference is what allows them to \emph{correct} the embedding signal rather than merely duplicate it.

Within those three constraints the ten features were chosen to span the biological evidence space with minimum redundancy. Features~1 and~2 (absolute target and gene counts) condition the remaining overlap features on graph coverage, preventing under-connected nodes from producing meaningless zero-overlap signals. Features~3 and~4 are the canonical direct-binding overlap score in both absolute and Jaccard-normalised form. Feature~5 promotes indirect but mechanistically meaningful overlap via shared pathway membership, which is important because many drug--disease pairs share pathways without sharing exact target genes. Feature~6 encodes prior-knowledge drug-class signal. Features~7 and~8 give a neighbour-of-a-neighbour view of chemical analogy and existing treatment evidence; features~9 and~10 do the same on the disease side. Together they cover the informal question ``can the molecule plausibly act on the relevant biology?'' through a handful of interpretable scalar features, each computable in $O(|T(c)| + |G(d)|)$ time per pair.

To minimise leakage, features~8 and~10 consult only the \emph{training} CtD edge set for the ``known treatment'' overlap, so test-time treatment edges are never revealed to the featuriser. Missing-data handling is explicit: if either $T(c)$ or $G(d)$ is empty (the compound has no recorded binding targets, or the disease has no recorded gene associations), the overlap and Jaccard features are set to zero and a binary indicator column is raised, so that the classifier learns to down-weight such under-informative pairs rather than treating a zero count as mechanistic dissimilarity.

On the 1{,}208-pair primary training set, Feature~3 (\texttt{shared\_targets}) has a mean of $\approx 0.31$ for positive CtD pairs versus $\approx 0.04$ for hard negatives --- an order-of-magnitude separation that pure embedding-based similarity does not capture. Feature~4 (\texttt{target\_overlap} Jaccard) is likewise roughly threefold larger for positives than for hard negatives on average. By contrast, the Abacavir~$\to$~ocular cancer pair --- the canonical score-validity inversion case --- has \texttt{shared\_targets}=0 and \texttt{target\_overlap}=0: there is no compound--disease gene overlap, only graph-structural proximity.

The MoA module adds 10 features to the classical feature vector (1177~$\to$~1187 dimensions, ${\sim}0.9\%$ relative increase) and is activated with \texttt{--use\_moa\_features}.

\subsection{Alleviating the score-validity inversion}

The score-validity inversion is a systematic bias, not a random error: embedding similarity rewards nodes that share many neighbours, and in a graph that encodes comorbidity and co-expression alongside direct treatment relations, that signal drifts away from mechanistic truth. The MoA features are designed to reintroduce the mechanistic axis into the scoring function in two complementary ways.

First, by \emph{down-weighting} pairs whose embedding score is high but whose mechanistic-evidence features are uniformly zero. In the stacking meta-learner of Section~\ref{sec:hybrid}, this manifests as a learned negative interaction between the embedding-based base-model probabilities and the MoA features: when MoA indicates no direct, no pathway, and no analogical evidence, the meta-learner reduces the final probability even though the base models are confident. Abacavir~$\to$~ocular cancer, which currently rises to the top of the ranking, is the prototypical candidate for this correction.

Second, by \emph{up-weighting} pairs whose embedding score is moderate but whose MoA features encode strong direct or analogical evidence. Losartan~$\to$~atherosclerosis --- angiotensin-receptor blockade, with documented anti-fibrotic and anti-inflammatory effects on vascular tissue --- scores only $\approx 0.52$ in the pure-embedding model despite seven Phase-4 clinical trials. Its MoA features, however, encode shared gene targets between Losartan and atherosclerosis-associated genes, overlap in cardiovascular pathways, and analogical treatment evidence from other angiotensin-receptor blockers and statins; together they provide a strong positive mechanistic signal that the meta-learner can use to promote the prediction.

Because the meta-learner is a linear classifier over a small number of base-model outputs augmented by ten interpretable features, the correction is transparent and can be audited feature-by-feature: a clinician can inspect which MoA features fired for any given prediction and verify that the post-correction score has a defensible mechanistic basis. This transparency is a first-order desideratum in clinical applications where regulatory and physician trust constrain usefulness. Quantifying the magnitude of the correction --- the change in Spearman correlation between model score and ClinicalTrials.gov trial count, and the rank change of Abacavir specifically --- is planned as the headline experiment of the v2 paper.

%% ==================================================================
\section{Experimental Setup}
%% ==================================================================

\subsection{Software and Hardware}

\begin{itemize}[leftmargin=*]
\item \textbf{Python} 3.9+; \textbf{PyKEEN} 1.10+ for KG embedding training; \textbf{Qiskit} 1.x and \textbf{Qiskit-Aer} for quantum circuits and statevector simulation; \textbf{scikit-learn} 1.4+ for classical models and stacking; \textbf{PyTorch} 2.x as PyKEEN backend.
\item \textbf{Quantum simulation:} Qiskit Aer CPU statevector simulator (primary); GPU-accelerated cuStateVec available via \texttt{--gpu}; IBM Quantum (Heron processor) for hardware runs.
\item \textbf{Memory:} ${\sim}$16 GB RAM for 16-qubit full kernel matrix; ${\sim}$8 GB for 12-qubit.
\item \textbf{Random seed:} 42 for all splits and training unless stated otherwise.
\end{itemize}

\subsection{Reproducing the Primary Result}\label{sec:exp_repro}

The following command reproduces the primary 0.7987 result:

\begin{verbatim}
python scripts/run_optimized_pipeline.py \
  --relation CtD \
  --full_graph_embeddings \
  --embedding_method RotatE \
  --embedding_dim 128 \
  --embedding_epochs 200 \
  --negative_sampling hard \
  --qml_dim 16 \
  --qml_feature_map Pauli \
  --qml_feature_map_reps 2 \
  --qsvc_C 0.1 \
  --qml_pre_pca_dim 24 \
  --run_ensemble \
  --ensemble_method stacking \
  --tune_classical \
  --fast_mode
\end{verbatim}

Expected QSVC kernel computation time: approximately 43 minutes on CPU. Results are written to \texttt{results/optimized\_results\_YYYYMMDD-HHMMSS.json}.

The \texttt{--fast\_mode} flag keeps the run practical: it uses reduced classical tuning (e.g., smaller tree ensembles) and skips optional heavy branches such as VQC comparisons in the same invocation, without changing the headline RotatE--Pauli--stacking configuration above.

%% ==================================================================
\section{Results}
%% ==================================================================

\subsection{Primary Results}

\begin{table*}[t]
\centering
\caption{Primary benchmark: test metrics for RotatE-128D with Pauli-16Q QSVC and stacking.\protect\\
\small\textit{Configuration: full-graph RotatE 128D (200 epochs), Pauli map reps${=}2$, PCA to 24D then 16Q, $C{=}0.1$, hard negatives 1:1, tuned classical models. Result file \texttt{optimized\_results\_20260216-100431.json}. QSVC kernel time 2,619~s (full $1208{\times}1208$ matrix).}}
\label{tab:primary}
\begin{tabular}{@{}lcccc@{}}
\toprule
\textbf{Model} & \textbf{Test PR-AUC} & \textbf{Test ROC-AUC} & \textbf{Test Accuracy} & \textbf{Type} \\
\midrule
\textbf{Ensemble-QC-stacking (Pauli)} & \textbf{0.7987} & 0.7456 & 0.5762 & Hybrid quantum-classical \\
RandomForest-Optimized & 0.7838 & 0.7319 & 0.5828 & Classical \\
ExtraTrees-Optimized & 0.7807 & 0.7301 & 0.6623 & Classical \\
Ensemble-QC-stacking (ZZ)$\dagger$ & 0.7408 & --- & 0.6490 & Hybrid quantum-classical \\
QSVC-Optimized (ZZ)$\dagger$ & 0.7216 & --- & 0.6556 & Quantum only \\
QSVC-Optimized (Pauli) & 0.6343 & 0.6313 & 0.5861 & Quantum only \\
\bottomrule
\end{tabular}

\smallskip
\begin{minipage}{\linewidth}
\small $\dagger$ZZ results from \texttt{results/optimized\_results\_20260216-091710.json} (same train/test split; QSVC genuine computation: 908 s). All Pauli-configuration rows are from the primary run.
\end{minipage}
\end{table*}

The stacking ensemble improves over the best classical model (RandomForest, 0.7838) by \textbf{+0.0149 PR-AUC} (absolute), or +1.49 percentage points. RF test precision is 1.000 with recall 0.166 --- indicating the model operates at very high precision on its confident predictions, consistent with the imbalanced nature of the underlying drug-disease relationship space.

\FloatBarrier
\subsection{Extended Results: 256D Embeddings and Optuna Tuning}

\begin{table*}[t]
\centering
\caption{Extended Configuration --- RotatE-256D, Pauli-12Q, Optuna. \textit{Configuration: RotatE 256D, 250 epochs, full-graph embeddings, Pauli feature map (reps=1), 12 qubits, no pre-PCA, QSVC C=0.676 (Optuna-tuned), hard negatives, stacking ensemble. Source: \texttt{results/optimized\_results\_20260323-134844.json}. Note: quantum kernel was computed once (${\sim}$99 s) on the first Optuna trial and reused across subsequent classical hyperparameter trials.}}
\label{tab:extended}
\begin{tabular}{@{}lcccc@{}}
\toprule
\textbf{Model} & \textbf{Test PR-AUC} & \textbf{Test ROC-AUC} & \textbf{Test Accuracy} & \textbf{Type} \\
\midrule
\textbf{Ensemble-QC-stacking (Pauli-256D)} & \textbf{0.8581} & 0.8245 & 0.7317 & Hybrid quantum-classical \\
RandomForest-Optimized & 0.8569 & 0.8231 & 0.7351 & Classical \\
ExtraTrees-Optimized & 0.8498 & --- & 0.7351 & Classical \\
QSVC-Optimized (Pauli-256D) & 0.7222 & 0.7272 & 0.6391 & Quantum only \\
\bottomrule
\end{tabular}
\end{table*}

The 256D configuration demonstrates that embedding dimensionality is the dominant performance driver. Upgrading from 128D to 256D RotatE embeddings lifts the classical RF baseline from 0.7838 to 0.8569 (+7.3 pp absolute) and the ensemble ceiling from 0.7987 to 0.8581 (+5.9 pp), while the structural pattern --- QSVC standalone below classical, ensemble above classical --- is preserved. The Optuna-tuned QSVC C=0.676 (vs.\ manual C=0.1) also contributes to the QSVC improvement (0.6343 $\to$ 0.7222).

\subsection{Feature Map Analysis: The Pauli Inversion Effect}
\label{sec:results_feature_map}

\begin{table*}[t]
\centering
\caption{Feature Map Ablation. \textit{RotatE-128D held constant; all other parameters identical.}}
\label{tab:ablation}
\begin{tabular}{@{}lccc@{}}
\toprule
\textbf{Feature Map} & \textbf{QSVC PR-AUC} & \textbf{Ensemble PR-AUC} & \textbf{$\Delta$ vs.\ RF} \\
\midrule
ZZ (reps=2, 908 s kernel) & \textbf{0.7216} & 0.7408 & $-0.043$ \\
\textbf{Pauli (reps=2, 2,619 s kernel)} & 0.6343 & \textbf{0.7987} & $+0.015$ \\
\bottomrule
\end{tabular}
\end{table*}

\begin{figure*}[t]
\centering
\includegraphics[width=0.92\textwidth]{../figures/fig2_pauli_zz.pdf}
\caption{Feature map selection reveals a counterintuitive quantum-classical tradeoff. Switching from ZZ to Pauli lowers standalone QSVC PR-AUC by 8.7 pp (0.7216 $\to$ 0.6343) while raising ensemble PR-AUC by 5.8 pp (0.7408 $\to$ 0.7987). The Pauli kernel generates predictions less correlated with classical tree model errors, providing greater complementary signal to the stacking meta-learner. Dashed line: best classical-only baseline (RandomForest, 0.7838).}
\label{fig:pauli_zz}
\end{figure*}

This is the central quantum finding of this work. Under an identical protocol, the Pauli map lowers standalone QSVC by 8.7~pp versus ZZ but raises the stacking ensemble by 5.8~pp. The ZZ stacking ensemble tops out at 0.7408---below RandomForest (0.7838)---whereas the Pauli stacking ensemble reaches 0.7987, above RF.

The explanation lies in prediction decorrelation. RandomForest and ExtraTrees learn axis-aligned partitions of the high-dimensional classical feature vector (1177 dimensions in the primary configuration). The ZZ kernel's structured Pauli-$Z$ interactions yield predictions moderately correlated with those tree partitions. The Pauli kernel---mixed Pauli interactions across qubit pairs---induces a different geometry; its errors align less with the forests'. The stacking meta-learner exploits that diversity: Pauli QSVC is individually weaker but \textit{more informative} as an additional base learner.

This result supports a broader principle: for quantum kernels in ensemble settings, the relevant optimization target is not standalone quantum accuracy but quantum-classical prediction \textit{decorrelation}.

\subsection{VQC Ablation}

Variational Quantum Circuit (VQC) results are reported for completeness. VQC is not included in the reported ensembles.

\begin{table*}[t]
\centering
\caption{VQC Ablation Results. \textit{Left: optimizer comparison (RealAmplitudes ansatz, 8 qubits, 50 iterations). Right: ansatz comparison (SPSA optimizer, 8 qubits, 50 iterations).}}
\label{tab:vqc_optimizer}
\begin{subtable}[t]{0.38\linewidth}
\centering
\caption{Optimizer comparison}
\label{tab:vqc_optimizer_sub}
\begin{tabular}{@{}lc@{}}
\toprule
\textbf{Optimizer} & \textbf{Test PR-AUC} \\
\midrule
SPSA & \textbf{0.5456} \\
COBYLA & 0.5086 \\
NFT & 0.4782 \\
\bottomrule
\end{tabular}
\end{subtable}
\hfill
\begin{subtable}[t]{0.58\linewidth}
\centering
\caption{Ansatz comparison}
\label{tab:vqc_ansatz}
\begin{tabular}{@{}lcc@{}}
\toprule
\textbf{Ansatz} & \textbf{Repetitions} & \textbf{Test PR-AUC} \\
\midrule
RealAmplitudes & 4 & \textbf{0.5474} \\
RealAmplitudes & 3 & 0.5342 \\
EfficientSU2 & 3 & 0.5173 \\
\bottomrule
\end{tabular}
\end{subtable}
\end{table*}

VQC performance in the current setup is near-random (approximately 0.50 baseline for a balanced dataset). SPSA outperforms gradient-free alternatives; deeper ansatzes (reps=4) marginally improve over shallower ones. The gap between VQC (${\sim}$0.55) and QSVC (0.63--0.72) suggests that the variational optimization landscape is challenging for these circuit depths and iteration budgets. VQC improvement via architecture search, warm-start initialization, and larger iteration budgets is deferred to future work.

\FloatBarrier
\subsection{Hardware Validation: IBM Quantum Heron}

IBM Quantum Heron hardware runs were conducted for a 16-qubit Pauli configuration. The hardware QSVC PR-AUC (${\approx}$0.634) is consistent with the statevector simulator result (0.6343), providing cross-validation that the simulator quantum kernel accurately reflects the hardware-executable computation. Full noise characterization and error mitigation analysis are planned for v2.

\FloatBarrier
%% ==================================================================
\section{Clinical Validation}
\label{sec:clinical}
%% ==================================================================

\subsection{Validation Methodology}

After identifying top-scoring novel (compound, disease) predictions --- pairs not present in the known CtD training edges --- we queried ClinicalTrials.gov for each predicted compound-disease combination using condition, intervention, and MeSH term filters. Trials in any phase (I--IV) and any status (completed, recruiting, active) were counted.

\subsection{Validation Results}

\begin{table*}[t]
\centering
\caption{ClinicalTrials.gov Validation of Top Novel Predictions}
\label{tab:clinical}
\small
\setlength{\tabcolsep}{3pt}
\begin{tabularx}{\textwidth}{@{}c >{\raggedright\arraybackslash}X c >{\raggedright\arraybackslash}p{0.14\textwidth} >{\raggedright\arraybackslash}p{0.12\textwidth} >{\raggedright\arraybackslash}X@{}}
\toprule
\textbf{Rank} & \textbf{Prediction} & \textbf{Score} & \textbf{Trials} & \textbf{Phase} & \textbf{Verdict} \\
\midrule
1 & Abacavir $\to$ Ocular Cancer & 0.793 & 0 & --- & Graph artifact; no clinical support. \\
2 & Ezetimibe $\to$ Gout & 0.693 & 0 direct; 4 anti-inflammatory & I--II & Novel plausible hypothesis. \\
3 & Ramipril $\to$ Stomach Cancer & 0.597 & 0 & --- & No clinical support. \\
4 & Losartan $\to$ Atherosclerosis & 0.528 & 7+ & Phase 4 & Strongly validated. \\
5 & Mitomycin $\to$ Liver Cancer & 0.525 & 7 (TACE) & Phase 2--3 & Strongly validated. \\
6 & Salmeterol $\to$ Liver Cancer & 0.520 & 0 & --- & No clinical support. \\
\bottomrule
\end{tabularx}
\end{table*}

\begin{figure*}[t]
\centering
\includegraphics[width=0.92\textwidth]{../figures/fig3_clinical.pdf}
\caption{Score-validity inversion: model prediction score vs.\ ClinicalTrials.gov registrations. The highest-scoring prediction (Abacavir$\to$ocular cancer, 0.793) has zero clinical trial support, while the most validated predictions (Losartan$\to$atherosclerosis, Mitomycin$\to$liver cancer, 7+ trials each) score 0.52--0.53. This inversion motivates the mechanism-of-action feature module.}
\label{fig:clinical}
\end{figure*}

\subsection{Score-Validity Inversion}

The highest-scoring prediction (Abacavir$\to$ocular cancer, 0.793) has zero clinical trial support, while the two most strongly validated predictions (Losartan$\to$atherosclerosis, Mitomycin$\to$liver cancer) score only 0.52--0.53. This \textit{score-validity inversion} is not a calibration failure in the usual sense --- the model correctly separates known CtD edges from non-edges --- but rather a \textit{false positive elevation} problem specific to multi-relational graph structure.

\textbf{Root cause.} Abacavir is a frontline HIV antiretroviral. HIV infection is associated with elevated incidence of ocular malignancies through immunosuppression pathways. Hetionet encodes these comorbidity relationships, and RotatE embeddings trained on the full graph absorb this structural proximity: Abacavir and ocular cancer entities end up close in embedding space through shared gene and disease neighbors, even though the pathway from HIV antiretroviral action to ocular cancer treatment is pharmacologically implausible.

This is the precise failure mode that the MoA feature module (Section~\ref{sec:moa}) is designed to address. Abacavir binds a small set of targets (reverse transcriptase, primarily) with zero overlap with the causal genes of ocular cancer. A \texttt{shared\_targets} feature value of 0 and a low \texttt{target\_overlap} Jaccard score would penalize this prediction regardless of its embedding-derived score.

\subsection{Novel Hypothesis: Ezetimibe for Gout}

Ezetimibe is a selective cholesterol absorption inhibitor acting at the NPC1L1 transporter. While no clinical trial directly investigates Ezetimibe for gout, four trials explore its anti-inflammatory properties in cardiovascular and metabolic contexts. The biological plausibility is grounded in emerging evidence linking lipid metabolism to urate transport: serum urate levels correlate with cholesterol, and NPC1L1 is expressed in proximal tubules relevant to urate reabsorption. This prediction represents a genuine, low-evidence hypothesis worthy of targeted investigation and serves as an example of the kind of mechanistically grounded novel candidate the pipeline is designed to surface.

%% ==================================================================
\section{Discussion}
%% ==================================================================

\subsection{Why the Ensemble Beats Classical-Only}

The stacking meta-learner treats QSVC as a base model whose predictions carry signal \textit{orthogonal} to RF and ET, even when its standalone accuracy is lower. RF and ET operate by axis-aligned splits on the 1177-dimensional classical feature vector; the Pauli quantum kernel encodes similarity via 16-qubit Hilbert-space inner products, with an inductive bias unlike coordinate-aligned trees. The meta-learner---logistic regression over base probabilities---weights QSVC as a ``second opinion from a different geometry,'' improving overall PR-AUC.

The quantitative evidence is clear: the Pauli kernel provides +5.8 pp ensemble gain over ZZ despite performing $-8.7$ pp worse in isolation. This decorrelation benefit is independent of the absolute QSVC accuracy level, suggesting that quantum kernel selection for ensembles should prioritize \textit{classical-quantum error independence} over quantum standalone performance.

\subsection{Embedding Quality as the Dominant Performance Driver}

Comparing Tables~\ref{tab:primary} and~\ref{tab:extended}: upgrading from 128D to 256D RotatE embeddings lifts every model --- RF (+7.3 pp), ET (+6.9 pp), QSVC (+8.8 pp), and ensemble (+5.9 pp). The full-graph training regime (all 2.25M Hetionet edges) means that entity representations for compounds and diseases are enriched by all 24 relation types. A compound's embedding encodes not just which diseases it treats, but which genes it binds, which pathways those genes participate in, what side effects it produces, and which pharmacologic class it belongs to. This multi-relational context is what makes RotatE embeddings rich enough to serve as meaningful quantum circuit inputs.

\subsection{Quantum Kernel Scaling and the \texorpdfstring{$O(n^2)$}{O(n-squared)} Constraint}

The fidelity quantum kernel requires $O(n^2)$ circuit evaluations: $n_{\text{train}}^2 / 2 = 1208^2/2 \approx 729\text{K}$ for the training kernel plus $n_{\text{train}} \times n_{\text{test}} \approx 182\text{K}$ for the test kernel, totaling approximately 911K evaluations. At 16 qubits with statevector simulation, this required 2,619 seconds. For the CbG relation (11,571 positive edges $\to$ ${\sim}$18,500 training pairs), naive full kernel computation would require approximately $18500^2/2 \approx 171\text{M}$ evaluations --- infeasible without Nystr\"{o}m approximation ($m$ landmarks: $O(nm)$ evaluations), hardware parallelism, or dataset subsampling. The 755-edge CtD relation sits at the practical boundary for exact quantum kernel computation on current simulators.

\subsection{Score-Validity Inversion and the MoA Solution}

The clinical validation reveals a structural weakness in pure embedding-based scoring: proximity in graph embedding space does not guarantee mechanistic plausibility. This is not unique to Hetionet --- any KG that encodes comorbidity, co-occurrence, or phenotypic similarity alongside direct treatment relations will produce false positive elevations for this reason.

The MoA feature module directly encodes the question the embedding cannot answer: \textit{does the compound interact with the biological machinery implicated in the disease?} Features 3--5 (\texttt{shared\_targets}, \texttt{target\_overlap}, \texttt{shared\_pathway\_genes}) provide direct mechanistic signal; features 8 and 10 (\texttt{similar\_compounds\_treat}, \texttt{similar\_diseases\_treated}) provide analogical treatment signal. Together, they give the model the vocabulary to distinguish ``structurally close in the graph'' from ``mechanistically plausible as a treatment.'' Empirical evaluation of MoA feature impact on the score-validity inversion is planned for v2.

\subsection{Limitations}

\textbf{Single relation.} All primary results use the CtD relation (755 positive edges). The pipeline is designed for multi-relational extension; CpD (Compound-palliates-Disease, 390 edges) and DrD (Disease-resembles-Disease, 543 edges) are natural immediate targets requiring no code changes beyond \texttt{--relation CpD}.

\textbf{Single seed.} Results in Tables~\ref{tab:primary} and~\ref{tab:extended} are single-seed runs (seed=42). No variance estimates are reported in v1. Multi-seed evaluation (3--5 seeds) is planned for v2 to provide confidence intervals and confirm result stability.

\textbf{Simulation-based quantum.} All primary results use CPU statevector simulation. IBM Quantum Heron hardware runs confirm comparable QSVC scores (${\sim}$0.634 hardware vs.\ 0.6343 simulator), but a complete hardware-equivalent benchmark with full error characterization is deferred to v2.

\textbf{MoA features not yet benchmarked.} The MoA module is fully implemented and integrated into the pipeline (\texttt{kg\_layer/moa\_features.py}, \texttt{kg\_layer/enhanced\_features.py}). Empirical PR-AUC results with \texttt{--use\_moa\_features} enabled will be reported in v2.

\textbf{No random or degree-heuristic baselines.} A random classifier (PR-AUC $\approx$ 0.50 for balanced data) and a degree-heuristic baseline (rank by compound degree in the training graph) are planned as explicit baseline rows for v2 Table~3.

\subsection{NISQ viability assessment: can IBM Heron actually confirm these gains?}

Before committing simulator-discovered improvements to hardware one must quantify what is achievable in principle and what is practically reachable on current noisy intermediate-scale quantum (NISQ) devices. Three sources of irreducible error bound the CtD ensemble from above: incomplete Hetionet curation (a fraction of ``hard negatives'' are likely unrecorded true treatments, implying a noise-limited PR-AUC ceiling of roughly $1 - 2\epsilon$ with $\epsilon$ estimated by re-labelling a random 5\,\% sample against DrugBank/DailyMed); \emph{coverage}-limited epistemic error (only 75 of the 137 diseases appear in training, so PR-AUC should be reported stratified by disease-coverage deciles); and hard-negative ambiguity near the decision boundary. Quantifying these three bounds gives a numerical target --- plausibly $0.88$--$0.92$ in the high-coverage stratum --- replacing the vague ``gap to 1.00'' with a scientifically defensible runway.

Against that ceiling, the most productive next steps are nearly all classical-side fixes that close the bottleneck between the 1{,}177-dimensional embedding vector and the 16-qubit quantum feature map: KTA-driven feature selection (replacing PCA with a kernel-target-alignment-maximising subset of semantically meaningful features); data re-uploading circuits that re-inject classical features with trainable rotations between layers; a trainable linear projection before the feature map, optimised to maximise KTA or minimise residual correlation with the best classical base learner; and an MoA-only ``late-fusion kernel'' stacked as an additional base learner whose signal is mechanistically separated from RotatE structural proximity. Mark's multi-Pauli conjecture --- that several Pauli feature maps at different circuit repetition counts act as weakly-correlated quantum base learners --- is testable in the same factorial design, conditioned on a KTA pre-screen that eliminates redundant (map, reps, qubits) tuples before any full kernel is computed. Two hard walls bound this axis: barren plateaus cause the kernel matrix to concentrate at reps~$\ge 3$ and $n > 16$ qubits, and ensemble diversity saturates once additional QSVC base learners carry residual correlation above roughly $0.7$ with those already in the ensemble; both should be monitored as explicit stop-criteria.

IBM Heron r2 (${\sim}156$ qubits, median two-qubit error $\sim 3$--$4\times 10^{-3}$, $T_{1}\approx 250\,\mu$s, $T_{2}\approx 150\,\mu$s, ECR-gate time ${\sim}68$\,ns) can execute a 16-qubit Pauli feature map at reps=2 within 1--2 coherence windows, and the hardware QSVC PR-AUC of $0.634$ reported in Section~7 confirms that the kernel signal survives the gate-error budget. With Nystr\"{o}m approximation at $m\approx 100$ landmarks, full ensemble evaluation reduces from ${\sim}1.46$M to ${\sim}150$K circuit evaluations, or roughly $40$--$60$ hours of queued hardware time --- tractable for a single targeted study but not for Optuna sweeps. What is not feasible on current Heron is the full factorial design: hardware noise preferentially flattens the Pauli kernel's off-diagonal structure, which is precisely the source of its decorrelation benefit in the stacking ensemble, so the simulator-observed Pauli-vs-ZZ gap ($+5.8$\,pp) may shrink to $<2$\,pp on hardware. The recommended protocol is therefore (i) run the full factorial design on CPU/GPU simulation with multi-seed variance bars; (ii) select the single configuration with the largest simulator-ensemble lift and the smallest expected hardware degradation --- from the current ablation, Pauli reps=2 at 12 qubits rather than 16; (iii) execute that configuration on Heron with Nystr\"{o}m$(m=100)$, zero-noise extrapolation, and M3 readout mitigation for 5 seeds; and (iv) require as the success criterion that the hardware ensemble PR-AUC exceeds the best noise-matched classical baseline by at least $0.02$ at $p<0.05$ over the 5 seeds.

The likelihood of this study succeeding on current hardware is moderate rather than high: a $+0.02$\,PR-AUC advantage over Random Forest is within reach of Heron with careful error mitigation but is not a safe bet, and a $+0.05$ advantage plausibly requires early fault-tolerant hardware on the 2028--2030 horizon of the IBM roadmap, conditional on logical-qubit yield. The scientific value of the experiment, however, does not depend on a positive outcome: even a negative result on current hardware --- paired with clear simulator evidence of the Pauli-kernel decorrelation benefit --- is a publishable contribution that quantifies the noise margin separating NISQ from utility in biomedical link prediction. The operational implication is that the most promising PR-AUC gains in the next 6--12 months will come from the classical-side fixes above, which Heron can then validate end-to-end, whereas Heron time itself should be budgeted as a validation run on the single best simulator configuration, not as an exploration tool.

%% ==================================================================
\section{Future Work}
%% ==================================================================

Future effort is organised around a single question: how can the primary ensemble PR-AUC of 0.7987 be pushed toward the irreducible ceiling of the CtD task without waiting for fault-tolerant quantum hardware? The central observation is that the largest near-term gains are recoverable on \emph{classical simulators} through better classical-to-quantum encoding and better use of the mechanism-of-action (MoA) signal, while the role of IBM Heron hardware is to \emph{validate}, not to \emph{explore}, the best simulator-selected configuration.

Table~\ref{tab:roadmap} summarises the planned simulator experiments in expected-PR-AUC-gain order, evaluated against the current primary configuration (RotatE-128D, Pauli-16Q, stacking ensemble, PR-AUC $= 0.7987$). A reasonable 6-month target is PR-AUC~$\approx 0.90 \pm 0.01$ on CtD, which sits comfortably inside the irreducible ceiling quantified in the preceding NISQ viability assessment.

\begin{table}[h]
\centering
\caption{Planned simulator experiments in expected-gain order.}
\label{tab:roadmap}
\small
\setlength{\tabcolsep}{6pt}
\renewcommand{\arraystretch}{1.25}
\begin{tabularx}{\columnwidth}{@{}>{\raggedright\arraybackslash}X>{\raggedleft\arraybackslash}p{0.28\columnwidth}>{\raggedright\arraybackslash}p{0.20\columnwidth}@{}}
\rowcolor{tabheader}
\textbf{Experiment} & \textbf{Expected gain} & \textbf{Cost} \\
\rowcolor{tabrowlight} Multi-seed CV of current best         & $\pm 0.015$ (variance estimate, no gain) & low \\
\rowcolor{tabrowdark}  KTA-driven feature selection          & $+0.02$ -- $0.04$                      & low \\
\rowcolor{tabrowlight} Data re-uploading + trained map       & $+0.03$ -- $0.05$                      & medium \\
\rowcolor{tabrowdark}  Multi-Pauli (reps $\in\{1,2,3\}$) stacked & $+0.01$ -- $0.03$                  & high (multiple kernels) \\
\rowcolor{tabrowlight} MoA-kernel late-fusion                & $+0.02$ -- $0.04$                      & low \\
\rowcolor{tabrowdark}  256-D embedding (already shown)       & $+0.06$ ($0.7987 \to 0.8581$)          & already done \\
\end{tabularx}
\end{table}

The following experiments and extensions are planned for v2 of this paper and subsequent work, grouped in time-horizon order and mapped directly onto the rows of Table~\ref{tab:roadmap}.

\vspace{0.6em}
\noindent\textbf{Immediate:}
\begin{itemize}[leftmargin=*]
\item \textbf{Multi-seed CV of the current best configuration} (Table~\ref{tab:roadmap}, row 1). Run the primary configuration (RotatE-128D, Pauli-16Q, reps=2, $C_{Q}=0.1$, hard negatives, stacking ensemble) across 5 independent seeds; report mean $\pm$ std PR-AUC and a 95\,\% confidence interval so that every subsequent gain can be tested for significance against a known variance baseline of roughly $\pm 0.015$. Add a degree-heuristic and a random classifier as explicit baseline rows on the same 5-seed grid.
\item \textbf{MoA-kernel late-fusion} (Table~\ref{tab:roadmap}, row 5). Activate the ten-feature MoA module via \texttt{--use\_moa\_features} on the primary CtD configuration and report PR-AUC before/after. In parallel, materialise the MoA features as a separate, small classical kernel stacked as an \emph{additional} base learner alongside RF, ET, LR, and QSVC; report the Spearman correlation between model score and ClinicalTrials.gov trial counts, and the rank change of the Abacavir $\to$ ocular cancer prediction under the score-validity inversion metric.
\item \textbf{KTA-driven feature selection} (Table~\ref{tab:roadmap}, row 2). Replace the 24-component PCA compression in the quantum path with the 24 raw-feature coordinates that maximise kernel-target alignment on the training labels. This preserves semantically meaningful axes (MoA counts, shared targets, Jaccard overlap) that PCA smears, at low compute cost; expected lift is $+0.02$ to $+0.04$ PR-AUC.
\end{itemize}

\vspace{0.6em}
\noindent\textbf{Medium-term:}
\begin{itemize}[leftmargin=*]
\item \textbf{Data re-uploading with a trained quantum-embedding map} (Table~\ref{tab:roadmap}, row 3). Replace the static linear projection into 16 qubits with a data re-uploading circuit in which the classical features are re-injected between trainable rotation layers; jointly train the pre-feature-map linear projection to maximise KTA (or to minimise residual correlation with the best classical base learner). Expected $+0.03$ to $+0.05$ PR-AUC at medium compute cost.
\item \textbf{Extension to the CpD and DrD relations.} Run the identical pipeline on Compound-palliates-Disease (390 edges) and Disease-resembles-Disease (543 edges) via \texttt{--relation CpD} / \texttt{--relation DrD}; report the first quantum KG results on these relations and test whether the Pauli-vs-ZZ ensemble-diversity gap transfers.
\item \textbf{Larger-scale relations with Nystr\"{o}m-approximated quantum kernels.} Extend to CbG (11{,}571 edges) and DaG (12{,}623 edges) using Nystr\"{o}m approximation at $m \approx 100$ landmarks --- the only currently tractable regime at these dataset sizes on statevector simulation.
\item \textbf{VQC ansatz search.} Larger circuit budgets (100+ iterations, reps $= 6$--$8$; RealAmplitudes, EfficientSU2, and TwoLocal ansatzes) to close the gap between the variational approach ($\sim 0.55$) and the QSVC baseline ($0.63$--$0.72$) observed in Table~6.
\end{itemize}

\vspace{0.6em}
\noindent\textbf{Longer-term:}
\begin{itemize}[leftmargin=*]
\item \textbf{Multi-Pauli stacked ensemble} (Table~\ref{tab:roadmap}, row 4). Add two or three additional Pauli feature maps with reps $\in \{1, 2, 3\}$ and optionally different qubit counts as independent quantum base learners, pre-screened for non-redundancy using kernel-target alignment and residual-correlation thresholds; stop adding configurations when residual correlation with the existing ensemble exceeds $\sim 0.7$ or when kernel entropy (variance of off-diagonal entries) plateaus. Expected $+0.01$ to $+0.03$ PR-AUC at high compute cost, with barren-plateau concentration at reps $\ge 3$ and $n > 16$ qubits as an explicit stop criterion.
\item \textbf{IBM Quantum Heron hardware validation.} Execute the single best simulator-selected configuration (likely Pauli reps=2 at 12 qubits rather than 16, per the recommended NISQ protocol) with Nystr\"{o}m$(m = 100)$, zero-noise extrapolation, and M3 readout mitigation across 5 seeds. Success criterion: hardware ensemble PR-AUC exceeds the best noise-matched classical baseline by at least $+0.02$ at $p < 0.05$.
\item \textbf{Variational quantum kernel learning.} Train the feature map jointly with the classifier as a successor to the fixed-kernel QSVC, once hardware maturity and optimiser stability permit end-to-end kernel training.
\item \textbf{DRKG-scale extension.} Apply the full pipeline to DRKG (4.4M edges, 97K entities) with Nystr\"{o}m-approximated quantum kernels as a stress test for scalability and multi-dataset generalisation.
\item \textbf{Multi-relational joint training.} Train on CtD + CpD + DrD simultaneously with a shared quantum feature space, building on the single-relation extensions above to test whether cross-relational signal sharing yields further ensemble diversity.
\end{itemize}

%% ==================================================================
\section{Conclusion}
%% ==================================================================

We have presented the first hybrid quantum-classical pipeline that bridges knowledge graph embeddings with quantum kernel classifiers for biomedical link prediction. On the Hetionet CtD task --- 755 known drug-disease treatments, 47,031 entities, 2.25 million edges across 24 relation types --- our system achieves:

\begin{itemize}[leftmargin=*]
\item \textbf{PR-AUC 0.7987} (primary: RotatE-128D, Pauli-16Q, genuine 2,619 s quantum kernel, +1.49 pp over best classical)
\item \textbf{PR-AUC 0.8581} (extended: RotatE-256D, Pauli-12Q, Optuna-tuned classical)
\end{itemize}

The central finding is a counterintuitive quantum-classical synergy: the Pauli feature map degrades standalone QSVC PR-AUC by 8.7 pp relative to ZZ but raises ensemble PR-AUC by 5.8 pp, because it generates predictions more decorrelated from classical tree model errors. This finding reframes the optimization target for quantum kernels in hybrid ensembles: maximize classical-quantum prediction independence, not standalone quantum accuracy.

Clinical validation against ClinicalTrials.gov reveals a score-validity inversion problem --- the highest-scoring prediction (Abacavir$\to$ocular cancer, 0.793) has zero clinical support while validated predictions score 0.52--0.53 --- and identifies Ezetimibe$\to$gout as a biologically plausible novel hypothesis. We introduce a ten-feature mechanism-of-action module encoding binding target overlap, pathway co-involvement, pharmacologic class, and treatment analogy to address this inversion.

This work establishes a new research direction at the intersection of quantum kernel methods and knowledge graph reasoning for biomedicine, with a reproducible open pipeline, verified experimental provenance, and a clear experimental roadmap toward multi-relation, multi-seed, and hardware-validated results.

%% ==================================================================
%% References
%% ==================================================================

\begin{thebibliography}{99}

\bibitem{himmelstein2017}
Himmelstein, D.S.\ et al.\ (2017).
Systematic integration of biomedical knowledge prioritizes drugs for repurposing.
\textit{eLife}, 6, e26726.
\url{https://doi.org/10.7554/eLife.26726}

\bibitem{mayers2023}
Mayers, M.\ et al.\ (2023; updated 2024).
Drug repurposing using consilience of knowledge graph completion methods.
\textit{bioRxiv}.
\url{https://doi.org/10.1101/2023.05.12.540594}

\bibitem{qkdti2025}
QKDTI (2025).
Quantum kernel-based drug-target interaction prediction.
\textit{Scientific Reports}.
\url{https://doi.org/10.1038/s41598-025-07303-z}

\bibitem{kruger2023}
Kruger, D.M.\ et al.\ (2023).
Quantum machine learning framework for virtual screening.
\textit{Machine Learning: Science and Technology}.
\url{https://doi.org/10.1088/2632-2153/acb900}

\bibitem{llmrag2025}
Deep learning-based drug repurposing using KG embeddings and GraphRAG (2025).
\textit{bioRxiv}.
\url{https://doi.org/10.64898/2025.12.08.693009}

\bibitem{quantum2026}
Large-scale quantum computing framework enhances drug discovery (2026).
\textit{bioRxiv}.
\url{https://doi.org/10.64898/2026.02.09.704961}

\bibitem{hybrid2024}
Hybrid classical-quantum pipeline for real-world drug discovery (2024).
\textit{bioRxiv}.
\url{https://doi.org/10.1101/2024.01.08.574600}

\bibitem{rotate2019}
Sun, Z.\ et al.\ (2019).
RotatE: Knowledge graph embedding by relational rotation in complex space.
\textit{ICLR 2019}.
\url{https://arxiv.org/abs/1902.10197}

\bibitem{pykeen2021}
Ali, M.\ et al.\ (2021).
PyKEEN 1.0: A Python library for training and evaluating knowledge graph embeddings.
\textit{Journal of Machine Learning Research}, 22(82), 1--6.
\url{https://jmlr.org/papers/v22/20-1531.html}

\bibitem{schuld2021}
Schuld, M.\ \& Petruccione, F.\ (2021).
\textit{Machine Learning with Quantum Computers}.
Springer.
\url{https://doi.org/10.1007/978-3-030-83098-4}

\bibitem{qiskit2023}
Qiskit contributors (2023).
Qiskit: An open-source framework for quantum computing.
\url{https://doi.org/10.5281/zenodo.2573505}

\bibitem{wolpert1992}
Wolpert, D.H.\ (1992).
Stacked generalization.
\textit{Neural Networks}, 5(2), 241--259.
\url{https://doi.org/10.1016/S0893-6080(05)80023-1}

\bibitem{havlicek2019}
Havlic\v{c}ek, V.\ et al.\ (2019).
Supervised learning with quantum-enhanced feature spaces.
\textit{Nature}, 567, 209--212.
\url{https://doi.org/10.1038/s41586-019-0980-2}

\bibitem{schuld2019}
Schuld, M.\ \& Killoran, N.\ (2019).
Quantum machine learning in feature Hilbert spaces.
\textit{Physical Review Letters}, 122, 040504.
\url{https://doi.org/10.1103/PhysRevLett.122.040504}

\bibitem{liu2021}
Liu, Y., Arunachalam, S., \& Temme, K.\ (2021).
A rigorous and robust quantum speed-up in supervised machine learning.
\textit{Nature Physics}, 17, 1013--1017.
\url{https://doi.org/10.1038/s41567-021-01287-z}

\bibitem{abbas2021}
Abbas, A.\ et al.\ (2021).
The power of quantum neural networks.
\textit{Nature Computational Science}, 1, 403--409.
\url{https://doi.org/10.1038/s43588-021-00084-1}

\bibitem{incudini2023}
Incudini, M.\ et al.\ (2023).
Resource saving via ensemble techniques for quantum neural networks.
\textit{Quantum Machine Intelligence}, 5, 23.
\url{https://doi.org/10.1007/s42484-023-00109-0}

\bibitem{optuna2019}
Akiba, T., Sano, S., Yanase, T., Ohta, T., \& Koyama, M.\ (2019).
Optuna: A next-generation hyperparameter optimization framework.
\textit{Proc.\ 25th ACM SIGKDD Int.\ Conf.\ on Knowledge Discovery \& Data Mining}, 2623--2631.
\url{https://doi.org/10.1145/3292500.3330701}

\end{thebibliography}

%% ==================================================================
%% Appendices
%% ==================================================================

\appendix

\section{Full Configuration: Primary Run}
\label{app:primary_config}

\textit{Complete configuration for \texttt{results/optimized\_results\_20260216-100431.json}.}

\begin{table*}[t]
\centering
\caption{Primary Run Configuration}
\label{tab:config_primary}
\begin{tabular}{@{}ll@{}}
\toprule
\textbf{Parameter} & \textbf{Value} \\
\midrule
\texttt{relation} & CtD \\
\texttt{embedding\_method} & RotatE \\
\texttt{embedding\_dim} & 128 \\
\texttt{embedding\_epochs} & 200 \\
\texttt{full\_graph\_embeddings} & True \\
\texttt{qml\_feature\_map} & Pauli \\
\texttt{qml\_feature\_map\_reps} & 2 \\
\texttt{qml\_dim} & 16 \\
\texttt{qml\_pre\_pca\_dim} & 24 \\
\texttt{qsvc\_C} & 0.1 \\
\texttt{qsvc\_nystrom\_m} & None (full kernel) \\
\texttt{negative\_sampling} & hard \\
\texttt{ensemble\_method} & stacking \\
\texttt{tune\_classical} & True \\
\texttt{fast\_mode} & True \\
\texttt{random\_state} & 42 \\
\texttt{cv\_folds} & 5 \\
QSVC kernel computation time & 2,618.6 s \\
Result file timestamp & 20260216-100431 \\
\bottomrule
\end{tabular}
\end{table*}

\section{Full Configuration: Extended Run}
\label{app:extended_config}

\textit{Complete configuration for \texttt{results/optimized\_results\_20260323-134844.json}.}

\begin{table*}[t]
\centering
\caption{Extended Run Configuration}
\label{tab:config_extended}
\begin{tabular}{@{}ll@{}}
\toprule
\textbf{Parameter} & \textbf{Value} \\
\midrule
\texttt{relation} & CtD \\
\texttt{embedding\_method} & RotatE \\
\texttt{embedding\_dim} & 256 \\
\texttt{embedding\_epochs} & 250 \\
\texttt{full\_graph\_embeddings} & True \\
\texttt{qml\_feature\_map} & Pauli \\
\texttt{qml\_feature\_map\_reps} & 1 \\
\texttt{qml\_dim} & 12 \\
\texttt{qml\_pre\_pca\_dim} & 0 (none) \\
\texttt{qsvc\_C} & 0.6756 (Optuna-tuned) \\
\texttt{qsvc\_nystrom\_m} & None \\
\texttt{negative\_sampling} & hard \\
\texttt{ensemble\_method} & stacking \\
\texttt{tune\_classical} & True \\
\texttt{fast\_mode} & True \\
\texttt{random\_state} & 42 \\
QSVC kernel computation (first trial) & ${\sim}$99 s \\
Kernel reuse across Optuna trials & Yes (83 trials, fixed quantum kernel) \\
Result file timestamp & 20260323-134844 \\
\bottomrule
\end{tabular}
\end{table*}

\end{document}
